% This is text associated with the open textbook "Open Electromagnetics"
% See immediately below for author, date
% This work is licensed under the Creative Commons Attribution-ShareAlike 4.0 International License. (CC BY SA 4.0) 
% To view a copy of the license, visit http://creativecommons.org/licenses/by-sa/4.0/

%ORB TITLE Permeability of Some Common Materials
%ORB AUTHOR 0        # S.W. Ellingson, Virginia Tech (ellingson@vt.edu)
%ORB DATE 20191128   # yyyymmdd
%ORB ABSTRACT Permeability of Some Commonly-Encountered Materials Materials

\label{m0136_Table_of_Permeabilities} % <-- Must match name of this file and module directory name.  Don't mess with this!
\index{permeability!of common materials|(}

The values below are relative permeability 
\index{permeability!relative}
$\mu_r \triangleq \mu/\mu_0$ for a few materials that are commonly encountered in electrical engineering applications, and for which $\mu_r$ is significantly different from $1$. These materials are predominantly ferromagnetic metals and (in the case of ferrites) materials containing significant ferromagnetic metal content.
Nearly all other materials exhibit $\mu_r$ that is not significantly different from that of free space.

The values presented here are gathered from a variety of references, including those indicated in ``Additional References'' at the end of this section.
Be aware that permeability may vary significantly with frequency; values given here are applicable to the frequency ranges for applications in which these materials are typically used.
Also be aware that materials exhibiting high permeability are also typically non-linear; that is, permeability depends on the magnitude of the magnetic field.
Again, values reported here are those applicable to applications in which these materials are typically used.  

{\bf Free Space} (vacuum): $\mu_r\triangleq 1$.

{\bf Iron} 
\index{iron}
(also referred to by the chemical notation ``Fe'') appears as a principal ingredient in many materials and alloys employed in electrical structures and devices.
Iron exhibits $\mu_r$ that is very high, but which decreases with decreasing purity.  
99.95\% pure iron exhibits $\mu_r\sim 200,000$.
This decreases to $\sim 5000$ at 99.8\% purity and is typically below 100 for purity less than 99\%. 

{\bf Steel} 
\index{steel}
is an iron alloy that comes in many forms, with a correspondingly broad range of permeabilities. 
\emph{Electrical steel}, commonly used in electrical machinery and transformers when high permeability is desired, exhibits $\mu_r\sim 4000$. % https://en.wikipedia.org/wiki/Electrical_steel
\emph{Stainless steel}, encompassing a broad range of alloys used in mechanical applications, exhibits $\mu_r$ in the range 750--1800. % https://en.wikipedia.org/wiki/Stainless_steel
\emph{Carbon steel}, including a broad class of alloys commonly used in structural applications,  exhibits $\mu_r$ on the order of 100. %https://en.wikipedia.org/wiki/Carbon_steel

{\bf Ferrites} 
\index{ferrite}
include a broad range of ceramic materials that are combined with iron and various combinations of other metals and are used as magnets and magnetic devices in various  electrical systems. Common ferrites exhibit $\mu_r$ in the range 16--640.

%--------------------------------

{\bf Additional Reading:}
\begin{itemize}
%\item Section~\ref{m0058_Magnetic_Materials} (``Magnetic Materials'')
\item \href{http://hbcponline.com}{\emph{CRC Handbook of Chemistry and Physics}}.
\item ``\href{https://www.microwaves101.com/encyclopedias/magnetic-materials}{Magnetic Materials}'' on \href{http://www.microwaves101.com}{microwaves101.com}.
\item ``\href{https://en.wikipedia.org/wiki/Permeability_(electromagnetism)}{Permeability (electromagnetism)}'' on Wikipedia.
\item ``\href{https://en.wikipedia.org/wiki/Iron}{Iron}'' on Wikipedia.
\item ``\href{https://en.wikipedia.org/wiki/Electrical_steel}{Electrical steel}'' on Wikipedia.
\item ``\href{https://en.wikipedia.org/wiki/Ferrite_(magnet)}{Ferrite (magnet)}'' on Wikipedia.
\end{itemize}

\index{permeability!of common materials|)}
